% SPAD Active/Passive Quenching Circuit
\documentclass[preview, border=10pt]{standalone}
\usepackage{tikz}
\usetikzlibrary{positioning, calc, arrows.meta, shapes, fit}
\usepackage{amsmath}
\usepackage{xcolor}
\begin{document}
\begin{tikzpicture}[
    scale=0.9,
    box/.style={draw, rounded corners=2pt, align=center, font=\small, minimum width=2.8cm, minimum height=0.9cm},
    sbox/.style={draw, rounded corners=2pt, align=center, font=\footnotesize, minimum width=2.2cm, minimum height=0.7cm},
    lbl/.style={font=\footnotesize},
]

% Title
\node[font=\bfseries] at (0, 2.5) {SPAD淬灭与复位电路 --- SiPM光子计数前端};

% === PASSIVE QUENCHING ===
\node[font=\bfseries, blue] at (-3.5, 1.5) {(a) 无源淬灭};
\node[sbox, fill=blue!5] (vapd) {$V_{APD}$\\{\footnotesize $>V_{BR}$}};
\node[sbox, fill=red!8, below=1.2cm of vapd] (rq) {淬灭电阻\\$R_Q$\\{\footnotesize 100k$\Omega$-1M$\Omega$}};
\node[sbox, fill=green!8, below=1.2cm of rq] (spad_p) {SPAD\\{\footnotesize Geiger模式}};

\draw[->] (vapd.south) -- (rq.north);
\draw[->] (rq.south) -- (spad_p.north);
\draw (spad_p.south) -- ++(0,-1.0) node[ground] {};

% Passive timing
\node[box, fill=blue!5, right=0.5cm of rq] (ptime) {
    \textbf{无源淬灭时序}\\[4pt]
    {\small 优点: 结构简单，面积小}\\[2pt]
    {\small 缺点: Dead Time 50-500ns}\\[2pt]
    {\small 光子计数率 < 10MHz}
};

\draw[->, dashed] (rq.east) -- (ptime.west);

% === ACTIVE QUENCHING ===
\node[font=\bfseries, blue] at (3.5, 1.5) {(b) 主动淬灭};

\node[sbox, fill=blue!5] (vapd2) {$V_{APD}$};
\node[box, fill=orange!10, below=1.0cm of vapd2, minimum width=3cm] (aq) {
    \textbf{主动淬灭反馈}\\[3pt]
    {\small 检测雪崩 → 拉低偏压}\\[2pt]
    {\small 雪崩停止 → 恢复偏压}};
\node[sbox, fill=green!8, below=1.0cm of aq] (spad_a) {SPAD\\{\footnotesize Geiger模式}};

\draw[->] (vapd2.south) -- (aq.north);
\draw[->] (aq.south) -- (spad_a.north);
\draw (spad_a.south) -- ++(0,-0.8) node[ground] {};

% Active timing
\node[box, fill=orange!5, right=0.5cm of aq] (atime) {
    \textbf{主动淬灭时序}\\[4pt]
    {\small 优点: Dead Time 2-10ns}\\[2pt]
    {\small 光子计数率 > 100MHz}\\[2pt]
    {\small 缺点: 电路复杂，功耗高}
};

\draw[->, dashed] (aq.east) -- (atime.west);

% SiPM Macro
\node[box, fill=cyan!8, below=3.5cm of spad_p, minimum width=14cm, xshift=2cm] (sipm) {
    \textbf{SiPM (硅光电倍增管) = N个SPAD微单元并联}\\[4pt]
    $I_\text{SiPM} = \dfrac{N_\text{fired}}{N_\text{micro}} \cdot G_\text{SPAD} \cdot q$\qquad
    $P_\text{det}=1-\exp\left(-\int_0^{T_g}[\eta_\text{PDE}\cdot\Phi_{ph}(t)+DCR]\,dt\right)$\\[4pt]
    {\footnotesize 线性区: $N_\text{fired}\ll N_\text{micro}$ \ $\Rightarrow$ \ $I_\text{SiPM}\propto\Phi_{ph}$ \quad|\quad
    饱和区: $N_\text{fired}\to N_\text{micro}$ \ $\Rightarrow$ \ 非线性响应}
};

\end{tikzpicture}
\end{document}
